% ch2_c (书页 73-84 / p088-p099)
%--B-TAIL--
%JOIN% 平直空间上总存在一个度规分量为常数的坐标系。

%FIG{2.21}: {二维球面上的线元。}
% ---- 书页 73 / p088 ----
但是，我们未必知道如何找到这样一个坐标系，而且空间偏离平直的方式有很多种；因此我们会想要一个对曲率更精确的刻画，这将在后面引入。

把$g_{\mu\nu}$化成它的\textbf{规范形式}（canonical form），便得到度规的一个有用刻画。在这种形式下，度规分量变为
\begin{equation*}
g_{\mu\nu}=\mathrm{diag}\,(-1,-1,\ldots,-1,+1,+1,\ldots,+1,0,0,\ldots,0),
\tag{2.46}
\end{equation*}
其中“diag”指具有所给元素的对角矩阵。度规的\textbf{号差}（signature）指正、负本征值两者的个数；例如，对 Minkowski 空间我们会说“号差为负正正正的度规”。如果本征值中有零，度规就是“退化的”（degenerate），逆度规将不存在；如果度规连续且非退化，它的号差在每一点都相同。我们要处理的将总是连续、非退化的度规。如果所有符号都是正的，这个度规就称为\textbf{欧几里得型}（Euclidean）或\textbf{黎曼型}（Riemannian）的（或者干脆叫\textbf{正定}（positive definite）的）；如果只有单独一个负号，则称为\textbf{洛伦兹型}（Lorentzian）或\textbf{伪黎曼型}（pseudo-Riemannian）的；而任何既有若干$+1$又有若干$-1$的度规都称为\textbf{不定}（indefinite）的。（于是，“欧几里得”一词有时意味着空间是平直的，有时又不是，但它始终意味着规范形式严格为正；这套术语令人遗憾，却是通行的标准。）广义相对论中感兴趣的时空都具有洛伦兹型度规。

我们还没有证明总可以把度规化成规范形式。事实上，在$M$的某一点$p$处总能做到这一点，但一般而言只在单独那一点处可行，在$p$的任何邻域内则不然。实际上我们还能做得比这稍微好一点；结果是，在任一点$p$都存在一个坐标系$x^{\hat\mu}$，使得$g_{\hat\mu\hat\nu}$在其中取其规范形式，而且一阶导数$\partial_{\hat\sigma}g_{\hat\mu\hat\nu}$全部为零（而二阶导数$\partial_{\hat\rho}\partial_{\hat\sigma}g_{\hat\mu\hat\nu}$则不可能全部变为零）：
% ---- 书页 74 / p089 ----
\begin{equation*}
g_{\hat\mu\hat\nu}(p)=\eta_{\hat\mu\hat\nu},\qquad \partial_{\hat\sigma}g_{\hat\mu\hat\nu}(p)=0.
\tag{2.47}
\end{equation*}
这样的坐标称为\textbf{局部惯性坐标系}（locally inertial coordinates），与之相联系的基矢量构成一个\textbf{局部洛伦兹标架}（local Lorentz frame）；当我们处于这些特殊坐标中时，常常给指标戴上帽号（hat）。注意，在局部惯性坐标系中，$p$点的度规在一阶近似下看上去就像平直空间的度规。这正是“时空的足够小区域看起来像平直的（Minkowski）空间”这一说法的严格含义。另外，在$M$的每一点上同时构造若干组\emph{基矢量}使得度规取其规范形式，这没有任何困难；问题在于，一般说来并不存在一个能从中导出这些基的\emph{坐标系}。这类基在附录 J 中有讨论。

关于如何构造局部惯性坐标系，我们将推迟到第 3 章再讨论。然而，对四维洛伦兹型度规这一具体情形勾画出其存在性的一个证明，是有益的。想法是考察度规的变换规律
\begin{equation*}
g_{\hat\mu\hat\nu}=\frac{\partial x^\mu}{\partial x^{\hat\mu}}\frac{\partial x^\nu}{\partial x^{\hat\nu}}\,g_{\mu\nu},
\tag{2.48}
\end{equation*}
并在待求的坐标$x^{\hat\mu}$下把两边作泰勒展开。旧坐标$x^\mu$的展开形如
\begin{align*}
x^\mu&=\left(\frac{\partial x^\mu}{\partial x^{\hat\mu}}\right)_p x^{\hat\mu}
+\frac{1}{2}\left(\frac{\partial^2 x^\mu}{\partial x^{\hat\mu_1}\partial x^{\hat\mu_2}}\right)_p x^{\hat\mu_1}x^{\hat\mu_2}\\
&\quad+\frac{1}{6}\left(\frac{\partial^3 x^\mu}{\partial x^{\hat\mu_1}\partial x^{\hat\mu_2}\partial x^{\hat\mu_3}}\right)_p x^{\hat\mu_1}x^{\hat\mu_2}x^{\hat\mu_3}+\cdots,
\tag{2.49}
\end{align*}
其余的展开沿同样的路线进行。[为简单起见，我们已令$x^\mu(p)=x^{\hat\mu}(p)=0$。]然后，用一套极其示意性的记号，(2.48)展开到二阶就是
\begin{align*}
&(\hat g)_p+\left(\hat\partial\hat g\right)_p\hat x+\left(\hat\partial\hat\partial\hat g\right)_p\hat x\hat x\\
&\quad=\left(\frac{\partial x}{\partial \hat x}\frac{\partial x}{\partial \hat x}g\right)_p
+\left(\frac{\partial x}{\partial \hat x}\frac{\partial^2 x}{\partial \hat x\,\partial \hat x}g+\frac{\partial x}{\partial \hat x}\frac{\partial x}{\partial \hat x}\hat\partial g\right)_p\hat x\\
&\quad\quad+\left(\frac{\partial x}{\partial \hat x}\frac{\partial^3 x}{\partial \hat x\,\partial \hat x\,\partial \hat x}g+\frac{\partial^2 x}{\partial \hat x\,\partial \hat x}\frac{\partial^2 x}{\partial \hat x\,\partial \hat x}g+\frac{\partial x}{\partial \hat x}\frac{\partial^2 x}{\partial \hat x\,\partial \hat x}\hat\partial g+\frac{\partial x}{\partial \hat x}\frac{\partial x}{\partial \hat x}\hat\partial\hat\partial g\right)_p\hat x\hat x.
\tag{2.50}
\end{align*}
我们可以令两边$\hat x$的同阶项彼此相等。于是，分量$g_{\hat\mu\hat\nu}(p)$——总共 10 个数（用来描述一个对称的二指标张量）——就由矩阵$\left(\partial x^\mu/\partial x^{\hat\mu}\right)_p$决定。这是一个$4\times4$
% ---- 书页 75 / p090 ----
矩阵，不带任何约束；因此我们可以自由选择 16 个数。显然，至少就“有足够的自由度”这一点而言，这些自由度足以把$g_{\hat\mu\hat\nu}(p)$的 10 个数化成规范形式。（实际上存在一些限制——如果你仔细走完整个程序，你会发现例如号差是无法改变的。）剩下的六个自由度恰好可以诠释为洛伦兹群的六个参数；我们知道这些变换保持规范形式不变。在一阶，我们有导数$\partial_{\hat\sigma}g_{\hat\mu\hat\nu}(p)$，四个导数乘十个分量，总计 40 个数。但观察 (2.50) 的右边我们看到，现在我们还有选取$\left(\partial^2 x^\mu/\partial x^{\hat\mu_1}\partial x^{\hat\mu_2}\right)_p$的额外自由。在这组数中，指标$\hat\mu_1$和$\hat\mu_2$有 10 种独立取法（它是对称的，因为偏导数相互对易），$\mu$有四种取法，总计 40 个自由度。这恰好是我们确定度规所有一阶导数所需的选取数目，因此我们可以把它们全部设为零。然而在二阶，我们要关心的是$\partial_{\hat\rho}\partial_{\hat\sigma}g_{\hat\mu\hat\nu}(p)$；它在$\hat\rho$与$\hat\sigma$以及$\hat\mu$与$\hat\nu$中对称，总计$10\times10=100$个数。我们作出额外选取的能力包含在$\left(\partial^3 x^\mu/\partial x^{\hat\mu_1}\partial x^{\hat\mu_2}\partial x^{\hat\mu_3}\right)_p$之中。它在三个下指标中对称，给出 20 种可能，再乘以上指标的 4，给我们 80 个自由度——比把度规的二阶导数全部设为零所需的少 20 个。所以事实上我们无法让二阶导数消失；因此对平直的偏离必须由代表度规张量场二阶导数的那 20 个自由度来量度。稍后我们将会看到这是怎么一回事，届时我们会用 Riemann 曲率张量来刻画曲率，而它将被证明在四维中有 20 个独立分量。

局部惯性坐标系有用得难以置信。最妙的是，它们的用处一般并不要求我们真的去做构造这种坐标的工作（尽管下一章我们会给出一个做法），而只要求我们知道它们确实存在。通常的窍门是：取一个有物理兴趣的问题，在局部惯性坐标系的框架下回答它，然后把答案表示成不依赖坐标的形式。举一个非常简单的例子：一位具有四速度$U^\mu$的观者，以及一枚以四速度$V^\mu$飞过的火箭。这位观者测得火箭的普通三速度是多少？在狭义相对论中答案很直接。在惯性坐标（全局地，而不仅局域地）中工作，让观者处于静止系，火箭沿$x$轴运动。于是观者的四速度是$U^\mu=(1,0,0,0)$，火箭的四速度是$V^\mu=(\gamma,\gamma v,0,0)$，其中$v$是三速度，而$\gamma=1/\sqrt{1-v^2}$，于是$v=\sqrt{1-\gamma^{-2}}$。由于（此刻）我们处在平直时空中，我们有
\begin{equation*}
\gamma=-\eta_{\hat\mu\hat\nu}U^{\hat\mu}V^{\hat\nu}=-U_{\hat\mu}V^{\hat\mu},
\tag{2.51}
\end{equation*}
因为$\eta_{00}=-1$。因此平直时空中的答案应当是
\begin{equation*}
v=\sqrt{1-\left(U_{\hat\mu}V^{\hat\mu}\right)^{-2}}.
\tag{2.52}
\end{equation*}
% ---- 书页 76 / p091 ----
现在我们可以回到弯曲时空了，在那里度规不再是平直的。但在进行测量的那一点上，我们可以随意使用局部惯性坐标系，此时$g_{\hat\mu\hat\nu}$的分量恰好就是$\eta_{\hat\mu\hat\nu}$的分量。所以在这个特定的坐标系中，(2.52)在弯曲时空中依然成立。但 (2.52) 是一个完全张量式的方程，它不在乎我们用的是什么坐标系；因此它以完全的一般性成立。这类做法将一再证明自己的价值。

\section{膨胀宇宙}

一个非平凡洛伦兹几何的简单例子由一个四维宇宙学时空提供，其度规为
\begin{equation*}
ds^2=-\mathrm{d}t^2+a^2(t)\left[\mathrm{d}x^2+\mathrm{d}y^2+\mathrm{d}z^2\right].
\tag{2.53}
\end{equation*}
这描写了一个宇宙，对它来说“某一固定时刻的空间”是一个平直的三维欧几里得空间，并且作为时间的函数在膨胀。保持在固定空间坐标$x^i$处不变的世界线称为共动的（comoving）；类似地，我们把随以固定空间坐标定义的边界一同膨胀的空间区域称为“共动体积”（comoving volume）。由于度规描写的是（距离）$^2$，在这个时空中，共动点之间的相对距离按$a(t)$增长；函数$a$称为尺度因子（scale factor）。这是 Robertson–Walker 度规的一个特殊情形，即空间切片在几何上平直的情形；还有另一些空间切片弯曲的情形（我们将在第 8 章讨论）。但我们此刻的兴趣不在于这个度规从何而来，而在于把它当作一个游乐场，用来演示我们已经发展出的某些想法。

尺度因子的典型解是幂律，
\begin{equation*}
a(t)=t^q,\qquad 0<q<1.
\tag{2.54}
\end{equation*}
实际上解五花八门，但上面这些是特别简单而又切题的几个。物质主导的平直宇宙满足$q=\frac{2}{3}$，而辐射主导的平直宇宙满足$q=\frac{1}{2}$。一个显著的特征是，当$t\to0$时尺度因子趋于零，度规的空间分量也随之趋于零。这是一个依赖于坐标的论断，原则上或许存在另一个坐标系使一切都显得有限；然而在这个例子中，$t=0$代表几何的一个真实奇点（即“大爆炸”），应当从流形中排除出去。因此$t$坐标的范围是
\begin{equation*}
0<t<\infty.
\tag{2.55}
\end{equation*}

我们的时空在$t=0$处走到尽头。

这个弯曲几何中的光锥由类光路径定义，即那些$ds^2=0$的路径。我们可以考虑使
% ---- 书页 77 / p092 ----
$y$和$z$保持不变的类光路径，来画出一幅时空图；于是
\begin{equation*}
0=-\mathrm{d}t^2+t^{2q}\,\mathrm{d}x^2,
\tag{2.56}
\end{equation*}
这意味着
\begin{equation*}
\frac{dx}{dt}=\pm t^{-q}.
\tag{2.57}
\end{equation*}
你也许会担心：我们前面才煞有介事地强调过$\mathrm{d}x^\mu$是基 1-形式而不是微分，转眼却马马虎虎地“用$\mathrm{d}t^2$作除法”从 (2.56) 得到了 (2.57)。实际情况要体面得多。我们实际所做的，是取由 (2.56) 定义的$(0,2)$型张量——它吃进两个矢量、返回一个实数——并把它作用在矢量$V=(dx^\mu/d\lambda)\partial_\mu$的两份拷贝上，$V$是曲线$x^\mu(\lambda)$的切矢量。只考虑$\mathrm{d}t^2$这一项作用在$V$上：
\begin{equation*}
\mathrm{d}t^2(V,V)\equiv(\mathrm{d}t\otimes\mathrm{d}t)(V,V)=\mathrm{d}t(V)\cdot\mathrm{d}t(V),
\tag{2.58}
\end{equation*}
其中记号$\mathrm{d}t(V)$指一个实数，我们按如下方式计算它：
\begin{align*}
\mathrm{d}t(V)&=\mathrm{d}t\left(\frac{dx^\mu}{d\lambda}\partial_\mu\right)\\
&=\frac{dx^\mu}{d\lambda}\,\mathrm{d}t\left(\partial_\mu\right)\\
&=\frac{dx^\mu}{d\lambda}\frac{\partial t}{\partial x^\mu}\\
&=\frac{dt}{d\lambda},
\tag{2.59}
\end{align*}
其中第三行用到了式 (2.25)。对$\mathrm{d}x^2$重复同样的步骤，我们发现 (2.56) 蕴含
\begin{equation*}
0=-\left(\frac{dt}{d\lambda}\right)^2+t^{2q}\left(\frac{dx}{d\lambda}\right)^2,
\tag{2.60}
\end{equation*}
由此经一维链式法则便得到 (2.57)：
\begin{equation*}
\frac{dx}{dt}=\frac{dx}{d\lambda}\frac{d\lambda}{dt}.
\tag{2.61}
\end{equation*}
教训应当很清楚：像 (2.56) 这样的表达式描写的是良定义的张量，但把基 1-形式当作单纯的“微分”来操作，确实能给出正确答案。（至少在大多数时候如此；把更正式的定义记在脑子里是个好主意。）

我们可以解 (2.57) 得到
\begin{equation*}
t=(1-q)^{1/(1-q)}\left(\pm x-x_0\right)^{1/(1-q)},
\tag{2.62}
\end{equation*}
% ---- 书页 78 / p093 ----
%FIG{2.22}: {平直 Robertson–Walker 宇宙（$a(t)\propto t^q$，$0<q<1$）的时空图。图中底部的虚线表示$t=0$处的奇点。由于光锥与奇点相切，两点的过去可以不重叠。}
其中$x_0$是积分常数。这些曲线定义了我们这个膨胀宇宙的光锥，如图 2.22 所绘。由于我们已假定$0<q<1$，光锥在$t=0$处与奇点相切。这个几何的一个关键特征是：两点的光锥在过去未必相交；这与 Minkowski 空间形成对照——在 Minkowski 空间中，任何两点的光锥在过去和未来都必定相交。我们说每个事件都定义了一个“视界”（horizon），在其之外存在对该事件处发生的事情不可能有任何影响的世界线。这是因为，既然没有任何东西能运动得比光快，每个点只能被那些或者位于其过去光锥之上、或者位于其内部的事件所影响（实际上，我们把过去光锥连同其内部合起来简称为一个事件的“过去”）。彼此处在对方视界之外的两个事件，称为“处于因果接触之外”（out of causal contact）。这些概念将在下一节以及第 4 章和第 8 章中得到更仔细的探讨。

\section{因果性}

许多物理问题都可以表述为初值问题（initial-value problem）：已知系统在某一时刻的状态，那么在稍后某时刻它的状态是什么？这类问题有确定答案，这件事归功于因果性（causality）——即这样一种观念：未来的事件可以被理解为初始条件加上物理定律所产生的后果。初值问题在 GR 中和在 Newton 物理学或狭义相对论中一样常见；然而，时空背景的动力学本性为初值表述的失效引入了新的可能方式。这里我们非常简略地介绍一些概念，它们用于理解因果性在 GR 中如何起作用。
% ---- 书页 79 / p094 ----
我们将考察在固定的背景时空上演化物质场的问题，而不是度规本身的演化。我们的指导原则是：任何信号都不能传播得比光速快；因此信息只会沿类时或类光的轨迹流动（不一定是测地线）。由于有时需要区分纯类时的路径和仅仅是非类空的路径，我们定义\textbf{因果曲线}（causal curve）为处处类时或类光的曲线。于是，给定流形$M$的任意子集$S$，我们定义$S$的\textbf{因果未来}（causal future），记作$J^+(S)$，为沿指向未来的因果曲线从$S$出发所能到达的点集；\textbf{编时未来}（chronological future）$I^+(S)$则是沿指向未来的类时曲线所能到达的点集。注意，零长度的曲线是因果的，但不是编时的；因此，点$p$永远在自己的因果未来$J^+(p)$之中，却不一定在自己的编时未来$I^+(p)$之中（尽管也可能在，下面会提到）。因果过去$J^-$与编时过去$I^-$可以类似地定义。

$M$的子集$S$如果没有任何两点能用类时曲线相连，就称为\textbf{非编时的}（achronal）；例如，Minkowski 时空中任何无边缘的类空超曲面都是非编时的。给定一个闭的非编时集$S$，我们定义$S$的\textbf{未来依赖域}（future domain of dependence），记作$D^+(S)$，为所有这样的点$p$的集合：经过$p$的\emph{每一条}向过去行进的不可延伸因果曲线都必须与$S$相交。（不可延伸（inextendible）只是说这条曲线永远延伸下去，不在某个有限点处终止；闭则是说该集合的补集是开集。）$S$自身的元素也是$D^+(S)$的元素。把未来换成过去，同样定义过去依赖域$D^-(S)$。一般而言，$M$中有些点会落在某个依赖域之内，有些则在其外；我们把$D^+(S)$的边界定义为\textbf{未来 Cauchy 视界}（future Cauchy horizon）$H^+(S)$，类似地把$D^-(S)$的边界定义为过去 Cauchy 视界$H^-(S)$。你可以自行确信它们都是类光曲面。依赖域和 Cauchy 视界绘于图 2.23 中，图中$S$取为非编时曲面$\Sigma$的一个连通子集。

%FIG{2.23}: {类空曲面$\Sigma$的一个连通子集$S$及其因果结构。$D^\pm(S)$表示$S$的未来/过去依赖域，$H^\pm(S)$表示未来/过去 Cauchy 视界。}
% ---- 书页 80 / p095 ----
这些定义的用处应当很明显：如果没有任何东西运动得比光快，信号就不可能传播到任一点$p$的光锥之外。因此，如果留在该光锥内部的每条曲线都必须与$S$相交，那么在$S$上给出的信息就应当足以预言$p$处的情形；也就是说，在$S$上给出的物质场的初始数据可以用来解出$p$处场的值。通过知道$S$上发生的事就能预言其情形的所有点的集合，是并集$D(S)=D^+(S)\cup D^-(S)$，简称\textbf{依赖域}（domain of dependence）。如果闭的非编时曲面$\Sigma$的依赖域$D(\Sigma)$是整个流形，就称$\Sigma$为\textbf{Cauchy 面}（Cauchy surface）；有了 Cauchy 面上给出的信息，我们就能预言整个时空中发生的事。如果一个时空拥有 Cauchy 面（它可能没有），就称它是\textbf{全局双曲}（globally hyperbolic）的。

%FIG{2.24}: {曲面$\Sigma$处处类空，但处在点$p$的过去光锥的过去；它的依赖域并不是整个时空。}
任何闭的、非编时的、且没有边缘的集合$\Sigma$都称为\textbf{部分 Cauchy 面}（partial Cauchy surface）。一个部分 Cauchy 面之所以未能成为真正的 Cauchy 面，可能归咎于它自己，也可能归咎于时空。一种可能性是我们恰好选了一张“坏”超曲面（尽管很难给出一个一般判据，说明超曲面何时在这种意义上是坏的）。考虑 Minkowski 空间以及一张无边缘的类空超曲面$\Sigma$，它停留在某个点的光锥的过去一侧，如图 2.24 所示。此时$\Sigma$是一张非编时曲面，但显然$D^+(\Sigma)$终止于光锥处，我们无法用$\Sigma$上的信息预言整个 Minkowski 空间中发生的事。当然，我们本可以挑别的曲面，使依赖域是整个流形，所以这件事并不怎么令我们担心。

Cauchy 视界出现的一种稍更不平凡的方式，是\textbf{闭合类时曲线}（closed timelike curves）的出现。在 Newton 物理学中，因果性靠时间的绝对观念冷酷无情地向前行进而得到强制。在狭义相对论中事情更加受限：你不仅必须在时间中前进，而且光速限制了你在空间中移动的迅捷程度（你必须待在自己的前向光锥之内）。在广义相对论中，你必须待在自己的前向光锥之内这一点仍然成立；然而它严格来说只是个局域的概念，因为从全局看，时空的曲率可能把光锥从一处“倾斜”到另一处。原则上，光锥可以被扭曲得足够厉害，使得一位观者能沿一条处处类时的前向路径运动，却与它自身在其“过去”的一点相交——这就是闭合类时曲线。

作为一个简单的例子，考虑一个带坐标$(t,x)$的二维几何，使坐标为$(t,x)$与$(t,x+1)$的点被等同起来。于是拓扑为$\mathbf{R}\times S^1$。我们取度规为
\begin{equation*}
ds^2=-\cos(\lambda)\,\mathrm{d}t^2-\sin(\lambda)\left[\mathrm{d}t\,\mathrm{d}x+\mathrm{d}x\,\mathrm{d}t\right]+\cos(\lambda)\,\mathrm{d}x^2,
\tag{2.63}
\end{equation*}
其中
\begin{equation*}
\lambda=\cot^{-1}t,
\tag{2.64}
\end{equation*}
% ---- 书页 81 / p096 ----
%FIG{2.25}: {一个带闭合类时曲线的柱状时空。光锥逐渐倾斜，使得曲面$\Sigma$的依赖域充满时空的下部，但当闭合类时曲线出现时，这个依赖域便告终结。}
它从$\lambda(t=-\infty)=0$过渡到$\lambda(t=\infty)=\pi$。这个度规并不代表广义相对论的什么著名的特解，它只是被杜撰出来充当闭合类时曲线的一个有趣例子；不过有一个众所周知的例子叫 Misner 空间（Misner space），性质与此类似。在 (2.63) 定义的时空中，光锥随着你走向未来而逐渐倾斜，如图 2.25 所示。当$t<0$时，光锥指向前方，因果性得以维持。然而一旦$t>0$，$x$就变成了类时坐标，于是可以沿一条绕$S^1$一圈回到自身的类时轨迹运动；这就是一条闭合类时曲线。如果我们先前指定的曲面$\Sigma$位于该点的过去，那么含闭合类时曲线的区域中的任何点都不在$\Sigma$的依赖域之内，因为闭合类时曲线自身不与$\Sigma$相交。因此在曲面$t=0$处必然存在一个 Cauchy 视界。这显然是比前一例更糟的问题，因为在这个时空中似乎不存在良定义的初值问题。

%FIG{2.26}: {$p$处的一个奇点把位于其未来的一切点，都从位于其过去的曲面$\Sigma$的依赖域中除去。}
最后一个例子由奇点的存在性提供：奇点是不在流形之中的点，尽管沿测地线走上有限距离就能到达它们。典型情形是曲率在某一点变为无穷大时出现奇点；一旦发生，这一点就不能再算是时空的一部分。这种情形可以导致 Cauchy 视界的出现，如图 2.26 所绘——点$p$处于某个奇点的未来，它就不可能位于该奇点过去某张超曲面的依赖域之内，因为从$p$出发会有一些曲线径直终止在奇点上。
% ---- 书页 82 / p097 ----
当我们试图从初始数据演化度规自身时，这些障碍同样会出现在 GR 的初值问题中。然而，它们的麻烦程度各不相同。挑到一张“坏”初始超曲面的可能性并不常出现，尤其因为大多数解都是（通过在整个时空中求解 Einstein 方程）全局地找到的。你必须小心的唯一情形是 Einstein 方程的数值求解：超曲面选得不好会导致数值困难，即使原则上完全解存在。闭合类时曲线似乎是 GR 极力避免的东西——肯定有包含它们的解，但从一般的初始数据演化通常不会产生它们。而奇点实际上不可避免。引力总是吸引的这一简单事实倾向于把物质拉到一起，使曲率增大，一般会导致某种奇点。显然我们必须学会与此共处，尽管存在这样的希望：一个良定义的量子引力理论会消除（或至少教会我们如何对付）经典 GR 的奇点。

\section{张量密度}

张量具有令人信服的美与简洁，但有些时候考虑非张量的对象是有用的。回忆第 1 章中我们引入的完全反对称的 Levi–Civita 符号，定义为
\begin{equation*}
\tilde\epsilon_{\mu_1\mu_2\cdots\mu_n}=
\begin{cases}
+1 & \text{若 $\mu_1\mu_2\cdots\mu_n$ 是 $01\cdots(n-1)$ 的偶置换（even permutation），}\\
-1 & \text{若 $\mu_1\mu_2\cdots\mu_n$ 是 $01\cdots(n-1)$ 的奇置换，}\\
0 & \text{其他情形。}
\end{cases}
\tag{2.65}
\end{equation*}
根据定义，Levi–Civita 符号在\emph{任何坐标系}中都具有上面规定的分量（至少在任何右手坐标系中是如此；交换手性会把$\tilde\epsilon_{\mu_1\mu_2\cdots\mu_n}$的分量乘上一个整体负号）。它之所以被称为“符号”，当然是因为它不是张量；它被定义为在坐标变换下不变。我们之所以只能在平直时空的惯性坐标中把它当作张量来处理，是因为洛伦兹变换反正不会改变其分量。它的行为可以与普通张量的行为联系起来，方法是首先注意到：给定任何$n\times n$矩阵$M^\mu{}_{\mu'}$，行列式$|M|$满足
\begin{equation*}
\tilde\epsilon_{\mu'_1\mu'_2\cdots\mu'_n}|M|=\tilde\epsilon_{\mu_1\mu_2\cdots\mu_n}\,M^{\mu_1}{}_{\mu'_1}M^{\mu_2}{}_{\mu'_2}\cdots M^{\mu_n}{}_{\mu'_n}.
\tag{2.66}
\end{equation*}
这不过是任何矩阵的行列式的一个紧凑表达式，与用代数余子式（cofactor）矩阵写出的通常公式完全等价。（你可以对$2\times2$或$3\times3$矩阵自行验证。）于是，令$M^\mu{}_{\mu'}=\partial x^\mu/\partial x^{\mu'}$，我们就得到
% ---- 书页 83 / p098 ----
\begin{equation*}
\tilde\epsilon_{\mu'_1\mu'_2\cdots\mu'_n}=\left|\frac{\partial x^{\mu'}}{\partial x^\mu}\right|\tilde\epsilon_{\mu_1\mu_2\cdots\mu_n}\,\frac{\partial x^{\mu_1}}{\partial x^{\mu'_1}}\frac{\partial x^{\mu_2}}{\partial x^{\mu'_2}}\cdots\frac{\partial x^{\mu_n}}{\partial x^{\mu'_n}},
\tag{2.67}
\end{equation*}
其中我们还利用了这样两个事实：矩阵$\partial x^{\mu'}/\partial x^\mu$是$\partial x^\mu/\partial x^{\mu'}$的逆，而逆矩阵的行列式是行列式的逆，$|M^{-1}|=|M|^{-1}$。于是，Levi–Civita 符号的变换方式与张量变换规律很接近，只是前面多出一个行列式。按这种方式变换的对象称为\textbf{张量密度}（tensor densities）。另一个例子由度规的行列式$g=|g_{\mu\nu}|$给出。容易验证，对 (2.48) 两边取行列式，就得到坐标变换下
\begin{equation*}
g(x^{\mu'})=\left|\frac{\partial x^{\mu'}}{\partial x^\mu}\right|^{-2}g(x^\mu).
\tag{2.68}
\end{equation*}
因此$g$也不是张量；它的变换方式与 Levi–Civita 符号类似，只是把雅可比行列式取到了$-2$次幂。雅可比行列式所取的幂称为张量密度的\textbf{权}（weight）；Levi–Civita 符号是权为 1 的密度，而$g$是权为$-2$的（标量）密度。

不过，我们不像喜欢张量那样喜欢张量密度。有一个简单的办法可以把密度变成诚实的张量——乘以$|g|^{w/2}$，其中$w$是密度的权（加绝对值符号是因为对洛伦兹型度规有$g<0$）。结果将按张量变换规律变换。因此，例如，我们可以把\textbf{Levi–Civita 张量}（Levi–Civita tensor）定义为
\begin{equation*}
\boxed{\ \epsilon_{\mu_1\mu_2\cdots\mu_n}=\sqrt{|g|}\,\tilde\epsilon_{\mu_1\mu_2\cdots\mu_n}.\ }
\tag{2.69}
\end{equation*}
由于这是一个真正的张量，我们可以升降指标等等。有时人们会定义一个带上指标的 Levi–Civita 符号版本$\tilde\epsilon^{\mu_1\mu_2\cdots\mu_n}$，其分量在数值上等于$\mathrm{sgn}(g)\tilde\epsilon_{\mu_1\mu_2\cdots\mu_n}$，其中$\mathrm{sgn}(g)$是度规行列式的符号。结果发现这是一个权为$-1$的密度，它与带上指标的张量（用$g^{\mu\nu}$把$\epsilon_{\mu_1\mu_2\cdots\mu_n}$的指标升起即得）之间的关系为
\begin{equation*}
\epsilon^{\mu_1\mu_2\cdots\mu_n}=\frac{1}{\sqrt{|g|}}\,\tilde\epsilon^{\mu_1\mu_2\cdots\mu_n}.
\tag{2.70}
\end{equation*}
你最终经常会做的事情之一，是把$\epsilon^{\mu_1\mu_2\cdots\mu_n}$的$p$个指标与$\epsilon_{\mu_1\mu_2\cdots\mu_n}$缩并；结果可以用 Kronecker δ 的反对称化乘积表示为
\begin{equation*}
\epsilon^{\mu_1\mu_2\cdots\mu_p\alpha_1\cdots\alpha_{n-p}}\epsilon_{\mu_1\mu_2\cdots\mu_p\beta_1\cdots\beta_{n-p}}=(-1)^s p!(n-p)!\,\delta^{[\alpha_1}_{\beta_1}\cdots\delta^{\alpha_{n-p}]}_{\beta_{n-p}},
\tag{2.71}
\end{equation*}
其中$s$是度规负本征值的个数（按我们的约定，洛伦兹号差下$s=1$）。最常见的例子是$p=n-1$，
% ---- 书页 84 / p099 ----
对此我们有
\begin{equation*}
\epsilon^{\mu_1\mu_2\cdots\mu_{n-1}\alpha}\epsilon_{\mu_1\mu_2\cdots\mu_{n-1}\beta}=(-1)^s (n-1)!\,\delta^\alpha_\beta.
\tag{2.72}
\end{equation*}

\section{微分形式}

现在让我们引入一类特殊的张量，称为\textbf{微分形式}（differential forms，或简称形式（form））。一个微分$p$形式不过是一个完全反对称的$(0,p)$型张量。于是，标量自动就是 0-形式，对偶矢量自动就是 1-形式（这就解释了此前的术语）。我们还有 4-形式$\epsilon_{\mu\nu\rho\sigma}$。所有$p$形式的空间记作$\Lambda^p$，流形$M$上所有$p$形式场的空间记作$\Lambda^p(M)$。一个半直截了当的组合练习揭示出，$n$维矢量空间上线性独立的$p$形式共有$n!/(p!(n-p)!)$个。于是在四维时空的一点上，有一个线性独立的 0-形式、四个 1-形式、六个 2-形式、四个 3-形式和一个 4-形式。$p>n$时不存在$p$形式，因为由反对称性，所有分量都自动为零。

为什么我们要在乎微分形式？这个问题不加些功夫难以回答，但基本想法是：形式既可以被微分，也可以被积分，而不需要任何额外几何结构的帮助。我们将对这两种操作各作一番简要的浏览。

给定一个$p$形式$A$和一个$q$形式$B$，取反对称化的张量积，我们可以构造出一个$(p+q)$形式，称为\textbf{楔积}（wedge product）$A\wedge B$：
\begin{equation*}
(A\wedge B)_{\mu_1\cdots\mu_{p+q}}=\frac{(p+q)!}{p!\,q!}\,A_{[\mu_1\cdots\mu_p}B_{\mu_{p+1}\cdots\mu_{p+q}]}.
\tag{2.73}
\end{equation*}
于是，例如两个 1-形式的楔积就是
\begin{equation*}
(A\wedge B)_{\mu\nu}=2A_{[\mu}B_{\nu]}=A_\mu B_\nu-A_\nu B_\mu.
\tag{2.74}
\end{equation*}
注意
\begin{equation*}
A\wedge B=(-1)^{pq}B\wedge A,
\tag{2.75}
\end{equation*}
所以只要你小心对待符号，就可以改变楔积的顺序。使用形式时我们可以随意省略指标，因为我们知道所有指标都在下边，而且张量是完全反对称的。

\textbf{外微分}（exterior derivative）$\mathrm{d}$使我们能够对$p$形式场求微分而得到$(p+1)$形式场。它定义为适当归一化的、反对称化的偏导数：
\begin{equation*}
(\mathrm{d}A)_{\mu_1\cdots\mu_{p+1}}=(p+1)\partial_{[\mu_1}A_{\mu_2\cdots\mu_{p+1}]}.
\tag{2.76}
\end{equation*}
最简单的例子是梯度，它就是 0-形式的外微分：
\begin{equation*}
(\mathrm{d}\phi)_\mu=\partial_\mu\phi.
\tag{2.77}
\end{equation*}
